\Titular%
{historia}%
{Relatos dunha vida extraordinaria}%
{Celia Álvarez Álvarez}%
{Memorias e leccións da a vida de Jane Goodall.}%

\begin{multicols}{2}

\subsection*{Introdución}

O pasado 1 de outubro faleceu a renombrada primatóloga e antropóloga
Jane Goodall, maiormente coñecida polo seu traballo de investigación
sobre a vida social dos chimpancés salvaxes do parque nacional Gombe
Stream, en Tanzania. En paralelo á investigación tamén levou a cabo
unha labor de activismo en materia de benestar animal e conservación.

\subsection*{Unha pequena científica}

\textit{Non é iso o que constitúe un pequeno científico? Curiosidade, facer
preguntas, non obter a resposta correcta, facer o intento de descubrir pola
conta dun, cometer un erro, non rendirse e aprender a ser paciente...}
%formateade esto como vexades, é unha cita dun vídeo

A historia de como Jane Goodall se converteu nunha científica comeza
na súa infancia. A pequena Jane adoraba pasar tempo fóra, xogando e
descubrindo a natureza. Jane vivía en Londres coa súa nai pero
adoraba os animais e o campo.

\indent Cando Jane tiña catro anos pasou un tempo nunha granxa do
campo inglés. A pequena nena integrouse nas labores do lugar
recolectando os ovos que poñían as galiñas e preguntando
insistentemente 'de onde sae o ovo?', 'se aquí está o ovo, onde está
o oco na galiña?'. Ninguén parecía interesado en explicarlle á
curiosa rapaza de onde viñan eses ovos, polo que un día foi ela mesma
a que se dispuxo a averigualo. Jane pasou catro horas dentro do
galiñeiro e saíu correndo cun brillo nos ollos para contarlle á súa
nai a marabillosa historia de como unha galiña pon un ovo. Durante
ese tempo ninguén sabía onde a nena se atopaba e a preocupación foi
tal que incluso chamaron á policía. Sen embargo, a nai de Jane non se
enfadou con ela nin a reprendeu por estar desaparecida durante esas
catro horas, senón que sentou a escoitar a historia que a súa filla
estaba tan emocionada por contarlle.

\indent Quizais se nese intre esa nai tivese pisado a curiosidade
científica desa nena, agora non estariamos falando de Jane Goodall.

\subsection*{Os comezos en Gombe}

Viaxar a África e vivir con animais era todo o que a xove Jane soñaba con
facer. Cando o tempo chegou, ir á universidade non era unha opción
económicamente asumíbel para a familia de Jane, polo que a rapaza conseguiu un
traballo de camareira. Gardaba cada propina para poder, algún día,  ir a
África.

\indent En 1957, durante a súa estancia na Colonia e Protectorado de Kenia
traballando como secretaria, Jane púxose en contacto co paleontólogo
Keniano-Británico Louis Leakey, co único desexo de poder conversar a cerca dos
animais. Louis Leakey estaba convencido de que o estudo dos grandes simios
podería proporcionar indicios sobre o comportamento dos primeiros homínidos e
buscaba un investigador de chimpancés. Jane Goodlall non tiña nin título nin
experiencia pero a Louis Leakey non pareceu importarlle. Todo o que estaba a
buscar era alguén de mente aberta, apaixoado polo coñecemento, namorado dos
animais e cunha paciencia monumental. Trala recaudación de fondos e unha
formación básica previa Goodall foi enviada en 1960 ao Parque Nacional de Gombe
Stream.

\indent Que unha moza solteira se instalase na África salvaxe non era
considerado seguro naqueles tempos, polo que Jane tivo que escoller un
acompañante. Foi a súa nai a que se presentou voluntaria, que aproveitou para
montar unha clínica e proporcionar mediciñas aos pacientes locais. Jane sempre
contou co apoio incondicional da súa figura materna: a súa nai nunca se opuxo
ás súas aspiracións. \textit{'Solamente es unha nena, non podes facer iso. Por
que non soñas con algo que poidas lograr?'} Son palabras que Jane Goodall
escoitou durante a súa vida, pero nunca ditas pola súa nai. \indent Durante os
inicios da súa labor investigadora, os chimpancés escapaban de Jane, estaban
horrorizados. Nunca antes viran un simio branco. Sentada, Jane observou por
semanas e semanas aqueles animais através dos seus binoculares, agardando que
dalgún xeito a comezasen a aceptar.

\subsection*{O Home, fabricante de ferramentas}

Tras un tempo Jane comezou a recibir visitas dun dos chimpancés, que atopara
plátanos polos que retornaba con frecuencia. O pelo branco no seu queixo foi
rasgo distintivo que permitía a Jane identificalo e a razón do seu apodo,
\textit{David Greybeard}. Cando Jane avistaba un grupo de chimpancés no bosque
recoñecía de inmediato a presencia de David, quen non fuxía e en cambio sentaba
tranquilamente. O resto de chimpancés, preparados para correr, miraban cara
David e debían pensar algo como 'ela non debe ser tan perigosa despois de
todo'.

\indent A medida que ía coñecendo aos membros do grupo, Jane púxolles nomes:
\textit{David Graybeard}, coa súa personalidade tranquila, solía estar
acompañado por \textit{Goliath}, o macho de maior rango naquel momento;
\textit{Míster McGregor}, un macho de avanzada idade cun carácter belixerante;
e despois estaba \textit{Flo}, de nariz bulboso e orellas esfiañadas, xunto coa
súa filla pequena Fifi. Cando Jane ollaba na mirada dun chimpancé vía unha
personalidade pensante e razoadora ollando de volta. Canto máis aprendía sobre
eles, máis se daba conta do semellantes que eran aos humanos en moitos
aspectos.

\indent Ao comezo da década dos anos sesenta era crido por moitos científicos
que soamente os humanos tiñan mentes. So os humanos eran capaces do pensamento
racional. Afortunadamente Jane Goodall non fora á universidade e e non sabía
esas cousas. Sentía como se estivera aprendendo a cerca de outros seres capaces
de sentir alegría, pena, medo e celos. Do mesmo xeito, durante moito tempo
tívose a crenza de que eramos as únicas criaturas na Terra que usaban e creaban
ferramentas.

Jane camiñaba pola longa e enmarañada vexetación do bosque cando de súpeto viu
unha forma escura encorvada sobre un montículo de cor dourada no que habitaba
un termiteiro. Apurando para coller os seus binoculares, Jane divisou as costas
dun chimpancé macho que non logrou identificar. O que puido ver en cambio foi
como este apañaba follas de herba e as introducía polo buraco do termiteiro.
Cando o animal se retirou Jane puido ver que se trataba de David Graybeard;
pero non foi ata que ela mesma colleu unha folla de herba e a introduciu polo
buraco que  comprendeu o que David estaba a facer. Ao retirar a folla decatouse
de como as termitas de aferraban a ela. \textit{O Home, fabricante de
ferramentas} era como se distinguía aos humanos doutros seres. E alí estaba
David Graybeard usando unha ferramenta. Uns días despois, Jane seguíu
observando aos chimpancés. Pescaban formigas, usaban follas como esponxas para
absorber auga, abrían noces cun martelo de madeira... O rudimentario comezo da
fabricación de ferramentas. Parecía que os humanos tiñan que ser redefinidos.

\indent A idea de que un animal puidese usar e crear unha ferramenta foi tan
impactante que a comunidade científica non deu crédito ás testemuñas de Jane
Goodall. Algúns tentaron desacreditar as observacións porque Jane era unha
muller xove e sen formación. Sen embargo o resultado foi a obtención dunha beca
da \textit{National Geographic Society} coa que Jane puido continuar os seus
estudos e coa que, a maiores, mandarían un fotógrafo para documentar o proceso.

\subsection*{O amor dunha nai}

Hugo fumaba. Sempre había cabichas no chan e a Jane non lle agradaba. El era un
perfeccionista e iso facía a Jane tolear. Sen embargo, era un home aposto e a
súa voz era calmada.

\indent A primeira noite que Hugo van Lawick pasou no campamento Jane escoitou
as historias a cerca dos filmes que creara na súa infancia e de como sempre
quixo fotografiar animais. Así que tiñan algo en común. Os dous amaban pasar
tempo ao aire libre, en contacto coa natureza, e os dous amaban o traballo que
estaban desempeñando. Así que remataron levándose moi ben.

\indent Para cando o tempo de Hugo en Gombe estaba chegando ao seu fin, este se
converteu nalguén moi importante para Jane. Ela sabía que o ía botar en falta.
Trala partida, Jane e Hugo intercambiaron varios telegramas.

\begin{center}
    \begin{minipage}{1\linewidth}
        \raggedright
        \texttt{WILL YOU MARRY ME STOP\\ LOVE STOP HUGO}
    \end{minipage}
    \begin{minipage}{1\linewidth}
        \raggedleft
        \texttt{YES STOP LOVE\\ JANE}
    \end{minipage}
    \begin{minipage}{1\linewidth}
        \raggedright
        \texttt{=DO YOU LIKE EMERALDS STOP\\ WHAT SIZE IS YOUR FINGER=\\ LOVE HUGO}
    \end{minipage}
    \begin{minipage}{1\linewidth}
        \raggedleft
        \texttt{LOVE EMERALDS STOP\\ LOVE YOU STOP\\ JANE}
    \end{minipage}
\end{center}

Hugo e Jane casaron, sen facer plans a futuro. Querían volver a Gombe e facer
filmes. Cando retornaron, foron recibidos cunha marabillosa sorpresa: Flo dera
a luz ao seu infante, Flint. E con Flint chegou a oportunidade de iniciar un
estudo que podería durar ata cincuenta anos. Era a primeira vez que se
observaría unha relación dunha nai chimpancé co seu infante de cerca e na
natureza.

\indent Flo era tódalas cousas que unha nai chimpancé debía ser: protectora
pero non sobreprotectora, agarimosa, lúdica e tolerante e usaba a distracción
en vez do castigo para ensinar ao seu infante. Fifi amosaba preocupación polo
seu pequeno irmán e Flo impedía delicadamente que se achegase a el. Isto non
durou moito pois pronto Flo permitíu a Fifi xogar con Flint e coidalo, e esta
se convertíu nunha grande axuda para a súa nai.

\indent O contacto cercano con Flo e a súa familia deixou unha fonda pegada no
espírito de Jane.

E Grub chegou. Jane e Hugo tiveron un fillo. Non se sentaran a falalo,
simplemente sucedeu naturalmente. Despois de que Flo dera a luz Jane pensou que
podería observar ao seu fillo. Jane planeaba levar a cabo un estudo do
desenvolvemento do seu cativo, tomando notas e capturando vídeo como fixera cos
chimpancés. Pero non funcionou. Jane quería estar presente en todo momento. E
como tódalas nais, Jane queríadarlle a Grub o mellor inicio posíbel na vida.
Como fonte de consello dispoñía da súa nai, pero tamén dispoñía de Flo. Non
cabía dúbida para Jane en canto ao moito que as súas observacións dos
chimpancés a axudaron a ser unha mellor nai, e en canto a súa experiencia de
converterse nunha nai a axudou a comprender mellor o comportamento maternal nos
chimpancés.

% better understand chimpanzee maternal Behavior it was not until grub came
% along for example that I began to understand the basic powerful instincts of
% Mother Love how much more easily I could Now understand the feelings of a
% chimpanzee mother furiously waved her arms and barked out threats to any who
% approached her infant too closely

\subsection*{Tempos escuros}

Hugo fora destinado ao Serengueti e Jane quedara en Gombe para continuar co seu
estudo. Novos e xoves investigadores foron destinados para acompañar a Jane.
Grub quedou en Gombe, coa súa nai. Pero tempos escuros chegaron: era horríbel,
un tras outro os chimpancés comezaban a mostrarse arrastrando as súas
extremidades. McGreggor non podía usar ámbalas pernas e máis un brazo. Podería
levarse a cabo o vaciñamento dos chimpancés, aínda que xa houbese algúns
afectados. Sen embargo, para McGreggor xa era demasiado tarde. Esperar a que a
natureza seguise o seu curso so prolongaría o seu sufrimento, tiña que ser
executado. Jane non podía ver a un animal sufrir así. \indent Despois do
incidente prohibiuse o contacto directo cos chimpancés e Gombe xa nunca foi o
mesmo.

Jane trataba de educar ao seu fillo e se sentía illada, botaba en falta a
compañía de Hugo. A vella Flo convertérase en avoa pois agora Fifi tiña un
infante do que coidar, dando paso a unha nova xeración na familia. Flint era xa
un adolescente e tiña unha idade na cal a maioría de machos xa pasaban tempo
lonxe da súa nai. Pero Flint seguía dependendo por completo de Flo e cada vez
que non obtía a atención da anciana nai comezaba unha intensa rabecha.

\indent Grub xa tiña idade para ir á escola polo que Jane tomou a decisión de
que comezaría o colexio en Inglaterra e viviría coa súa avoa. Jane sentiu unha
grande pena cando se despediu de Grub e voltou para Gombe. No nadal e en
primaveira Jane voltaría ao Reino Unido e nos veráns el a visitaría en
Tanzania.

E Flo faleceu. Cruzaba un regato de grande corrente de auga. A escena resultaba
pacífica, coma se o seu corazón parase de palpitar de súpeto. Flint sentaba na
beira do regato preto do corpo de Flo e se achegaba pedindo que o reconfortase
como sempre fixo durante a súa vida. A depresión de Flint empeoraba, deixou de
comer e adoitaba estar illado. Finalmente enfermou, parecía que sen a súa nai
xa non tivese ánimos para vivir. Cerca de tres semanas despois da morte de Flo,
Flint faleceu tamén.

Trala morte de Flo a comunidade de chimpancés que tan ben coñecía Jane comezou
a dividirse. Un grupo de chimpancés desprazouse cara ao sur do territorio,
separándose e perdendo o seu dereito de ser tratados como membros da
comunidade. A calma dese mundo ídilico foi perturbada e comezou unha guerra. O
grupo que se movera cara ao sur quedou completamente aniquilado.

\indent Jane non imaxinaba que os chimpancés fosen capaces de semellante
brutalidade pois sempre pensara que a guerra era un comportamento puramente
humano. Custoulle asimilar que o lado escuro e malvado da natureza humana
estaba profundamente arraigado nos xenes herdados dos nosos antepasados.

O matrimonio de Jane e Hugo viuse afectado pola distancia. O traballo de Hugo
agora estaba no Serengueti e quedarse en Gombe o afastaría dos seus obxectivos.
Era dificultoso continuar o matrimonio soamente por Grub e as discusións
entrambos tornáronse frecuentes. Hugo desexaba que Jane o acompañase no
Serengueti, pero o estudo en Gombe significaba moito para ela.

\indent A separación resultou moi triste para todos, especialmente para Grub.
Pero a experiencia de Jane no bosque lle outorgara perspectiva. No bosque a
morte non está escondida e rodea a un en todo momento. O ciclo permanente da
vida: os chimpancés nacen, fanse maiores, enferman e logo morren; e sempre
están os membros xoves para perpetuar a vida da especie.

\subsection*{Cada un de nós}

O noso intelecto creceu enormemente en complexidade dende as nosas orixes fai
dous millóns de anos e solamente os humanos fomos capaces de desenvolver unha
linguaxe oral sofisticada. Por primeira vez na evolución unha especie é capaz
de ensinar aos seus membros máis xoves a cerca de obxectos e sucesos non
presentes, de pasar a sabedoría adquirida a través de éxitos e erros do pasado.
Coa linguaxe podemos preguntar como ningún outro ser vivo pode facer, esas
preguntas a cerca de quenes somos e por que estamos aquí. E este intelecto
altamente desenvolto significa que temos unha responsabilidade cara ás outras
formas de vida do noso planeta, cuxa existencia vese ameazada polo
comportamento irreflexivo da nosa especie.

A vida de Jane Goodall nesa época era mellor que calquera cousa que ela nunca
tivera imaxinado; a pesares de todo seguía pasando o seu tempo no campo, estaba
escribindo un libro e á investigación en Gombe chegaban novos estudantes. Pero
sabía que os chimapncés en África estaban desaparecendo e deuse conta de que
debía aumentar a concienciación ao respecto.

Os humanos non estamos exentos da extinción e debemos tomar conciencia do
importante que é preservar o medio ambiente e o resto de especies que o
habitan. O que cada un de nós fai ten un impacto. Cada día que vivimos temos un
impacto no planeta e podemos escoller que clase de impacto queremos ter. Cada
un de nós é unha única persoa e os problemas globais nos poden provocar
desesperación. Pero podemos comezar pensando na nosa comunidade, que podes
facer na túa comunidade? Podes comezar facendo algo que a mellore, animar a
outros a que che axuden. Verás que podes facer unha diferenza. E iso se sinte
ben. Quererás facer máis e inspirarás a máis persoas. E quizais atrévaste a
pensar globalmente.

\indent Estas son algunhas das ideas tralo activismo de Jane Goodall, ao cal se
adicou dende entón e ata o seu último día.

\textit{Cando boto a vista atrás á miña vida, parece que tiven unha sorte
extraordinaria, aínda que, como sempre di a miña nai, Vanne, a sorte foi só
unha parte da historia, ela sempre creu que o éxito vén a través da
determinación e o traballo duro e que a culpa non está nas nosas estrelas,
senón en nós mesmos, que somos subordinados. Certamente creo que iso é certo,
pero aínda que traballei duro toda a miña vida, debo admitir que as estrelas
tamén parecen ter xogado o seu papel.}
%formateade esta cita como vexades, non creo convinte poñer a fonte tampouco, é
%da peli-documental pero como a meirande parte do contido

\subsection*{Nota da autora}

Se o lector chegou ao final deste artigo, quizais se pregunte cal é o sentido
que ten incluilo nunha revista universitaria de Física. Con esta pequena nota
pretendo transmitir a miña humilde crenza de que a Física non está illada
doutras disciplinas e de que como aspirantes a ou xa físicos, científicos e por
riba de todo, humanos (pois tamén contemplo o caso de que o lector sexa alleo á
Ciencia e de que se atope lendo esta revista por interese ou casualidade); non
debemos descoidar a nosa humanidade. Os valores e as motivacións que Jane
Goodall expresou durante a súa vida e que deixou co seu legado lograron
espertar o meu recordo pola miña humanidade, que ás veces se perde entre o
barullo do diario.

\indent Segundo a RAG, a definición de \textit{humanidade} é a seguinte:
'\textit{comprensión, compaixón e benevolencia cara aos demais}'.   Considero a
historia de Jane Goodall coma un exemplo de amor cara aos animais, cara á
Natureza e cara á vida. O interese de coñecer e descubrir as verdades da
Natureza é innato aos seres humanos, polo que penso que de certo modo esta
curiosidade tamén é capaz de espertar amor en nós por todo aquelo que nos rodea
como elementos dun conxunto ou dun \textit{todo}. Ás veces resúltame bo lembrar
o motivo polo que estou estudando física e polo cal sempre debo tentar
respectar tódolos elementos do Universo, do cal eu tamén formo parte.

\nocite{youtube_lastInterview}
\nocite{youtube_insideLook}
\nocite{youtube_discussingBehavior}
\nocite{goodall_wikipedia}
\nocite{bbc_horizon}

\printbibliography

\end{multicols}
